\section{Dogma and the deposit of faith}

A dogma is a divinely revealed truth proposed by the Church for belief as
revealed. Its formal motive is God's authority revealing, not the historical
plausibility of a committee or the private brilliance of a pope. The Church's
proposal supplies the authoritative ecclesial identification of the object;
it does not cause the truth to become revealed. Vatican I, \emph{Dei Filius},
chapter 3, names both forms of infallible proposal: solemn judgment and the
ordinary and universal Magisterium.

This excludes two reductions. Dogma is not confined to propositions bearing a
particular verbal charm such as ``we define.'' Intention and object must be
manifest, but infallible teaching also occurs through the bishops' universal
and concordant proposal in communion with Peter. Conversely, emphatic papal
language, canonization, liturgical celebration, disciplinary severity, or
frequent repetition does not by itself transform every associated statement
into dogma.

A definition is delimited. The Church may define that Christ is one divine
person in two natures without defining every philosophical account of
personhood; she may define Mary's bodily Assumption without defining whether
Mary died. Fidelity includes the definition's deliberate silences. The
contrary error often begins by making a dogma assert more than its competent
act asserts and then rejecting the exaggeration.

